\chapter{REVISIÓN CRÍTICA DE LA LITERATURA}

%Colocar base de datos que se usa segun el accuracy por ejemplo que sale,
La investigación en Control Tolerante a Fallas (\textit{Fault-Tolerant Control}, FTC) para vehículos aéreos no tripulados tipo cuadricóptero ha evolucionado de manera significativa, dando lugar al desarrollo de diversas estrategias orientadas a mejorar la confiabilidad y seguridad de estos sistemas frente a fallas en sensores, actuadores y perturbaciones externas \cite{freddi2019,fdd2021survey,quadcopter2024}. Como resultado, se han propuesto enfoques basados en modelado matemático, estimación de estados, aprendizaje profundo y aprendizaje por refuerzo, cada uno con diferentes capacidades y limitaciones para la detección, diagnóstico y compensación de fallas \cite{fdd2021survey,quadcopter2024,ozcan2024,liu2024,kim2025}. En este contexto, la presente revisión crítica de la literatura analiza las principales contribuciones reportadas en el área con el propósito de identificar sus fortalezas, limitaciones y las brechas de investigación que motivan el desarrollo de la propuesta planteada en esta tesis. Para ello, los trabajos se organizan en cuatro ejes temáticos: (i) métodos basados en modelos y observadores, (ii) enfoques de aprendizaje profundo para la detección de fallas, (iii) estrategias de Control Tolerante a Fallas basadas en aprendizaje por refuerzo y (iv) enfoques híbridos de adaptación en línea.


\section{Métodos Basados en Modelos y Observadores de Estado}


\noindent Los enfoques clásicos para la detección y compensación de fallas en cuadricópteros se fundamentan en el conocimiento del modelo dinámico del sistema, donde observadores o estimadores generan señales residuales cuya desviación respecto al comportamiento nominal indica la presencia de una anomalía.

\noindent Freddi, Longhi y Monteriù~\cite{freddi2019} propusieron un sistema integrado que combina observadores dinámicos con un controlador pasivo basado en realimentación, detectando fallas en actuadores mediante análisis de residuales y empleando un banco de observadores reducidos para el diagnóstico de fallas en sensores. Si bien el esquema fue validado en MATLAB/Simulink y constituye una referencia fundacional del área, la detección queda supeditada a la precisión del modelo linealizado, lo que restringe su capacidad ante dinámicas no lineales y perturbaciones aerodinámicas externas. En una dirección complementaria centrada exclusivamente en sensores, Patan y Caliskan~\cite{patan2022} compararon el filtro de Kalman extendido (EKF), el unscented (UKF) y el de cubatura (CKF) para la estimación de estado bajo fallas en sensores, concluyendo que el CKF minimiza el error de estimación. No obstante, ninguno de estos enfoques integra un mecanismo de compensación en lazo cerrado que reconfigure el controlador tras la detección, ni contempla fallas simultáneas en actuadores y sensores.

Amoozgar, Chamseddine y Zhang~\cite{amoozgar2013} aportaron una validación sobre hardware real mediante un TSKF que estima el vector de efectividad de los actuadores en la plataforma Quanser Qball-X4, tanto para fallas en un único motor como para fallas simultáneas en los cuatro actuadores. Este trabajo destaca por ser uno de los pocos del área con pruebas en tiempo real sobre un cuadricóptero físico; sin embargo, el TSKF asume que el proceso estocástico de las fallas es conocido a priori, lo que compromete su rendimiento ante fallas progresivas, y no incorpora un módulo que reconfigure automáticamente el controlador tras la detección.

\noindent Para superar esta brecha entre estimación y compensación activa, Du et al.\cite{du2023} propusieron combinar el Control de Rechazo Activo de Perturbaciones (ADRC) con un filtro de Kalman de dos etapas (TSKF), donde el TSKF estima el factor de pérdida de efectividad del actuador y un Observador de Estado Extendido (ESO, \textit{...}) dentro del lazo ADRC compensa el error residual de estimación. Esta integración mejora la estabilidad y el rechazo de perturbaciones respecto a controladores PID clásicos en simulación, aunque la validación permanece exclusivamente en MATLAB/Simulink sin considerar fallas en sensores. Complementando este enfoque desde el área de los UAVs de ala fija, Bekhiti et al.~\cite{bekhiti2025} emplearon redes neuronales de función de base radial (RBF-NN) con inversión dinámica no lineal, incorporando un observador no lineal para la estimación de fallas en sensores y la red RBF-NN para la compensación adaptativa de fallas en actuadores bajo condiciones de hielo. Aunque este trabajo representa el enfoque más completo de la sección al tratar fallas simultáneas en sensores y actuadores, su arquitectura está diseñada para ala fija y no es transferible directamente a cuadricópteros.


\section{Detección de Fallas Basada en Aprendizaje Profundo}

Con el avance del aprendizaje profundo han surgido enfoques orientados a la detección temprana de fallas mediante el análisis de patrones secuenciales en datos de vuelo, prescindiendo de un modelo matemático explícito del sistema.

Ozcan y Perkgoz~\cite{ozcan2024} propusieron una arquitectura que combina un autoencoder LSTM para reducción de dimensionalidad con un clasificador BiLSTM para análisis bidireccional de las secuencias temporales, alcanzando una precisión del 79.48\% en la predicción de fallas hasta 30 segundos antes de que ocurra un problema grave, empleando datos de vuelo reales recolectados con giroscopio, acelerómetro, magnetómetro, GPS y sensores de corriente a bordo. Complementando esta perspectiva desde un enfoque semi-supervisado, Li et al.~\cite{li2021} aplicaron un modelo LSTM-OCSVM que detecta desviaciones anómalas en señales multivariadas de sensores sin necesidad de ejemplos etiquetados de falla durante el entrenamiento, lo que resulta especialmente útil en escenarios donde los datos de falla son escasos. Ambos trabajos comparten una limitación estructural: se concentra en la detección sin incorporar ningún esquema de compensación activo que permita al UAV mantener la estabilidad una vez identificada la anomalía.

En contraste, Van Schijndel y Sun~\cite{schijndel2021} optaron por una arquitectura de bajo costo computacional basada en un estimador de filtro de Kalman que calcula un factor de efectividad estocástico para cada actuador usando únicamente mediciones de RPM, giróscopo y acelerómetro a bordo, integrando además un test de hipótesis para identificar el actuador fallido. A diferencia de los enfoques anteriores, este algoritmo fue validado en vuelo real en tiempo real y acoplado con un sistema de FTC activo, lo que lo convierte en el trabajo más cercano a un despliegue embebido real dentro de esta sección. Su principal restricción es que asume parámetros aerodinámicos nominales conocidos y no contempla fallas en sensores ni perturbaciones externas que puedan enmascarar la señal de falla.


\section{Control por Modo Deslizante con Observadores No Lineales}

\noindent Se combina el control por modo deslizante (SMC) con observadores no lineales para lograr compensación simultánea de fallas en actuadores y perturbaciones externas, sin depender de un modelo linealizado del sistema.

Wang, Ma y Gao~\cite{wang2022} presentaron un controlador de modo deslizante terminal no singular rápido (NFTSMC) con un observador de perturbaciones que estima de forma conjunta el efecto de las fallas en actuadores y las perturbaciones de viento, garantizando convergencia en tiempo finito. Chen et al.~\cite{chen2023} extendieron este paradigma incorporando una superficie de deslizamiento integral y un observador de perturbaciones que compensan simultáneamente fallas en actuadores y perturbaciones externas, reportando seguimiento preciso de trayectoria en simulación. Ambos trabajos son validados en MATLAB/Simulink, no contemplan fallas en sensores y presentan el riesgo de \textit{chattering} ante perturbaciones de alta frecuencia no modeladas.

Desde una perspectiva de validación más cercana al hardware, Ahmadi et al.~\cite{ahmadi2023} propusieron un controlador FTC activo con SMC basado en un observador no lineal que estima las fallas de motor como incertidumbre paramétrica interna, comparándose favorablemente con controladores INDI (\textit{Incremental Nonlinear Dynamic Inversion}) y MINDI (\textit{Modified Incremental Nonlinear Dynamic Inversion}) tanto en simulación no lineal como en experimentos Hardware-in-the-Loop (HIL). Aunque la validación HIL representa un avance respecto a la simulación pura, el sistema se enfoca exclusivamente en fallas en motores sin abordar fallas en sensores ni perturbaciones aerodinámicas externas, y la validación HIL no equivale a pruebas en vuelo real. Gao et al.~\cite{gao2022}, por su parte, adoptaron un enfoque híbrido modelo-datos donde un ESO estima la perturbación total y un clasificador Deep Forest determina el tipo y magnitud de la falla en actuadores con bajo costo computacional. Si bien este trabajo ofrece diagnóstico preciso y eficiente, no integra compensación que reconfigure el controlador tras el diagnóstico, ni contempla fallas en sensores ni perturbaciones externas simultáneas.


\section{Control Tolerante a Fallas Mediante Aprendizaje por Refuerzo}


El aprendizaje por refuerzo (RL,... ) ha evolucionado para el FTC en cuadricópteros, dado que permite aprender políticas de control flexibles sin requerir un modelo explícito del sistema y puede adaptarse a condiciones dinámicas no vistas durante el entrenamiento.


Liu, Yuan, Gao y Zhang~\cite{liu2024} propusieron un marco que combina una política PPO con una ley de alcance de modo deslizante terminal y un observador no lineal de perturbaciones, donde la política RL genera señales compensatorias sobre el controlador base PID. A diferencia de los enfoques SMC de la sección anterior, este trabajo fue validado tanto en simulación como en experimentos en el mundo real, demostrando robustez ante pérdidas de efectividad de hasta el 100\% en un único rotor. Sohège, Quiñones-Grueiro y Provan~\cite{sohege2021} abordaron el problema desde otra perspectiva, combinando un controlador supervisado RL con PIDs de bajo nivel e incorporando garantías formales de estabilidad que evitan que la política viole restricciones de seguridad. Sin embargo, ambos trabajos comparten la limitación de abordar únicamente fallas en actuadores: el primero puede no generalizar bien ante fallas simultáneas en múltiples motores, y el segundo carece de módulo de detección y diagnóstico de fallas, sin identificar el tipo ni la magnitud de la anomalía.

Desde una perspectiva centrada en fallas en sensores, Fei, Tu, Xu y Deng~\cite{fei2020} desarrollaron Learn-to-Recover, que combina un controlador base robusto con una política RL entrenada para recuperarse de desviaciones extremas causadas por señales de sensores falsificadas bajo ataques ciberfísicos. Aunque el trabajo demuestra la viabilidad del RL para recuperación ante fallas en sensores, aborda un escenario muy específico sin contemplar fallas en actuadores, degradación gradual de componentes ni perturbaciones aerodinámicas, y la transferibilidad a hardware real no fue evaluada.

Integrando las fortalezas de estos enfoques, Kim, Lee, Bang y Bae~\cite{kim2025} propusieron un marco híbrido que en la Fase 1 entrena conjuntamente un codificador de factores ambientales (MLP de 2 capas, dimensión latente de 8) y una política RL mediante PPO, y en la Fase 2 entrena un módulo Transformer (2 capas de codificador, dimensión del modelo 128, 2 cabezas de atención) para inferir el vector latente a partir del historial de estado-acción de 1 segundo sin acceder a parámetros privilegiados del sistema. Validado en el simulador PyBullet, el método alcanzó una tasa de éxito del 95\% y un RMSE posicional de 0.129m bajo pérdidas de efectividad del 30\% en un motor, superando a la línea base RL-RMA (86\% de éxito, RMSE~=~0.153~m) y al controlador PID estándar (53\% de éxito), con generalización a configuraciones de UAV no vistas. No obstante, el sistema no aborda fallas en sensores y la validación permanece exclusivamente en simulación sin evaluación en plataformas embebidas de bajo costo.

\section{Estrategias Híbridas de Adaptación en Tiempo Real}


Un método emergente en el FTC para sistemas autónomos es el diseño de módulos de adaptación en tiempo real capaces de ajustar la política de control sin reentrenamiento, enfrentando dinámicas no vistas durante el entrenamiento.

El paradigma de este método es el algoritmo RMA (\textit{Rapid Motor Adaptation}) de Kumar, Fu, Pathak y Malik~\cite{kumar2021}, concebido originalmente para robots cuadrúpedos. En la Fase 1 se entrena una política base (MLP de 3 capas) con acceso a factores ambientales privilegiados codificados en un vector extrínsico de 8 dimensiones; en la Fase 2, un módulo de adaptación CNN-1D estima ese vector a partir del historial de 50 pasos de estado-acción. En despliegue, ambos módulos corren de forma asíncrona (10~Hz y 100~Hz respectivamente) sin ajuste fino, demostrando robustez en terrenos complejos no vistos sobre el robot A1 de Unitree. Este paradigma fue adoptado directamente por Kim et al.~\cite{kim2025} en el contexto de FTC para UAVs, sustituyendo el módulo CNN-1D por un Transformer para mejorar la inferencia ante cambios abruptos de estado. La extensión de RMA a dinámicas aéreas no es trivial, dado que los cuadricópteros presentan mayor inestabilidad y menor tolerancia a las oscilaciones durante la adaptación que los robots terrestres.

Por otro lado, Giral, Gómez, Vinuesa y Le-Clainche~\cite{giral2024} propusieron un Transformer que mapea directamente valores de referencia de bucle externo a comandos de control en UAVs de ala fija, eliminando los controladores internos y las capas de detección de fallas mediante destilación de conocimiento maestro-alumno. Zhang et al.~\cite{zhang2023}, por su parte, presentaron un controlador de hovering único para cuadricópteros con configuraciones físicas muy distintas (masa e inercia hasta 10 veces diferentes), validado tanto en simulación como en hardware real, que infiere parámetros físicos a partir de datos de vuelo recientes sin reentrenamiento. Aunque este trabajo es el más cercano al despliegue real dentro de esta sección, está orientado a la adaptación ante variaciones físicas del cuadricóptero y no a la tolerancia a fallas explícita en actuadores o sensores. Finalmente, Hsu et al.~\cite{hsu2024} propusieron ReForMa, donde un agente adversario adaptativo genera durante el entrenamiento las perturbaciones externas más difíciles para el controlador, logrando mayor robustez ante viento y variaciones de carga en simulación PyBullet que los métodos de \textit{domain randomization} estándar. Sin embargo, el marco no aborda fallas en actuadores ni en sensores y su entrenamiento adversarial puede ser inestable ante ajustes inadecuados de hiperparámetros.



La revisión de los trabajos presentados permite identificar patrones y brechas estructurales en el estado del arte. Desde el punto de vista de la arquitectura, los métodos basados en observadores~\cite{freddi2019,patan2022,du2023,bekhiti2025,amoozgar2013,wang2022,chen2023,ahmadi2023,gao2022} ofrecen eficiencia computacional y estimación precisa del estado bajo condiciones nominales, pero dependen de la exactitud del modelo y presentan capacidad limitada para capturar dinámicas no lineales complejas ante fallas simultáneas. Los enfoques de aprendizaje profundo~\cite{ozcan2024,li2021} identifican patrones anómalos sin modelo explícito, pero requieren grandes volúmenes de datos etiquetados. Los métodos de aprendizaje por refuerzo~\cite{liu2024,sohege2021,fei2020,kim2025} aprenden políticas flexibles adaptables a condiciones dinámicas, aunque la mayoría presenta brechas de transferencia simulación-realidad. Los enfoques híbridos de adaptación en tiempo real~\cite{kumar2021,giral2024,zhang2023,hsu2024,kim2025} representan el paradigma más reciente, combinando políticas aprendidas con módulos de adaptación temporal que permiten ajuste en tiempo real sin reentrenamiento.

\section{Consideraciones finales}

Respecto a la cobertura de fallas, la mayoría de los trabajos aborda de forma independiente fallas en actuadores~\cite{freddi2019,wang2022,chen2023,ahmadi2023,liu2024,sohege2021,kim2025,amoozgar2013} o fallas en sensores~\cite{patan2022,bekhiti2025,fei2020}, sin integrar su interacción simultánea con perturbaciones externas de origen aerodinámico. Esta fragmentación metodológica evidencia una desconexión entre la detección de fallas, la estimación de estados y la compensación en control. En cuanto a la viabilidad computacional, únicamente el trabajo de Van Schijndel y Sun~\cite{schijndel2021} valida explícitamente su solución en tiempo real sobre hardware de vuelo embebido. El resto realiza evaluaciones en simulación (MATLAB/Simulink, PyBullet, RaiSim) o HIL~\cite{ahmadi2023}, sin abordar los requisitos de cómputo de plataformas embebidas de bajo costo.

En consecuencia, existe una brecha técnica relevante: se requiere un enfoque integrado que combine estimación de estados y aprendizaje temporal ligero para la detección y compensación conjunta de fallas en sensores y actuadores, considerando perturbaciones dinámicas externas, en tiempo real y con eficiencia computacional adecuada para plataformas embebidas. 